\section{Related Work}
\label{sec:related}

\subsection{Video Anomaly Detection: from Reconstruction to Weak Supervision and Open Vocabulary}

Early VAD methods adopt reconstruction or future-frame prediction, using the criterion that normal samples reconstruct well while anomalies do not. They require no anomaly samples but struggle to separate ``hard to reconstruct'' from ``genuinely abnormal'', and lack robustness to semantic variation in complex scenes. To reduce annotation cost, the weakly supervised setting became dominant: training with video-level labels only and pushing that supervision down to clips through multiple-instance learning. \citeauthor{yun2025missiongnn}~et al.\ combine a hierarchical multimodal graph neural network with mission-specific knowledge graph generation, improving structured reasoning for weakly supervised anomaly recognition~\cite{yun2025missiongnn}. Nevertheless, weakly supervised methods remain tied to predefined anomaly categories and generalise poorly to unseen ones.

Open-set and open-vocabulary directions attempt to break this limit. \citeauthor{wu2024openvocab}~et al.\ propose open-vocabulary VAD so that the model can recognise anomaly categories unseen during training~\cite{wu2024openvocab}; \citeauthor{kim2026dtvad}~et al.\ further improve temporal localisation in the open-vocabulary setting with a dual-branch aligned temporal-aware framework~\cite{kim2026dtvad}. Recent surveys systematically review a decade of method evolution and evaluation protocols~\cite{abdalla2025vad10years,pravalika2026surveillance} and point out that evaluation dimensions beyond detection scores have long been missing. A recent survey on violence detection reaches a similar conclusion~\cite{houssein2026violence}.

\subsection{Vision-Language Large Models for Anomaly Detection and Understanding}

Vision-language models map video and text into a shared semantic space, providing open-vocabulary criteria for anomaly detection. \citeauthor{zanella2024clip}~et al.\ were among the first to exploit the CLIP latent space for anomaly recognition, demonstrating the potential of zero-shot transfer~\cite{zanella2024clip}. Subsequent work extends this line in several directions: \citeauthor{biswas2025mmvad}~et al.\ build a cross-domain vision-language VAD model with contrastive learning and scale-adaptive frame segmentation~\cite{biswas2025mmvad}; \citeauthor{kang2025clipalign}~et al.\ use adaptive CLIP semantic alignment for efficient weakly supervised detection~\cite{kang2025clipalign}; \citeauthor{ahn2026anyanomaly}~et al.\ propose zero-shot customisable LVLM anomaly detection~\cite{ahn2026anyanomaly}; \citeauthor{kebour2025zeroshot}~et al.\ apply a zero-shot vision-language model to event detection in smart surveillance~\cite{kebour2025zeroshot}; \citeauthor{zou2026prompting}~et al.\ unlock VLM anomaly detection with fine-grained prompting~\cite{zou2026prompting}, while \citeauthor{ignatosyan2026prompt}~et al.\ specifically study and mitigate prompt sensitivity~\cite{ignatosyan2026prompt}. Work aimed at understanding rather than detection alone is also emerging: \citeauthor{zhou2026targetvau}~et al.\ propose multimodal anomaly-aware reasoning for target behaviour understanding~\cite{zhou2026targetvau}, \citeauthor{li2026aerial}~et al.\ exploit multimodal large language models for semantic-aware aerial VAD~\cite{li2026aerial}, and \citeauthor{yang2025monitor}~et al.\ perform online VAD via instruction following~\cite{yang2025monitor}.

Together these works establish that large models are usable for VAD, but their evaluation remains dominated by detection metrics and generally assumes well-behaved inputs; the score depression and threshold drift caused by domain shift are rarely quantified.

\subsection{Training-Free Paradigms and Retrieval-Augmented Reasoning}

The training-free paradigm reuses pre-trained large models without in-domain retraining, which lowers deployment cost and eases cross-domain transfer. \citeauthor{zanella2024trainingfree}~et al.\ demonstrate training-free anomaly localisation with a large language model~\cite{zanella2024trainingfree}; \citeauthor{pei2026drvad}~et al.\ improve the stability of training-free criteria via definition-guided reasoning~\cite{pei2026drvad}; \citeauthor{sun2026rag4vad}~et al.\ bring retrieval-augmented generation into the explanation stage so that explanations can rely on external knowledge~\cite{sun2026rag4vad}. Retrieval augmentation resonates with the knowledge-graph line: \citeauthor{wang2025holotrace}~et al.\ build a bidirectional causal knowledge graph for edge--cloud collaborative detection~\cite{wang2025holotrace}.

The key risk of the training-free paradigm is its implicit assumption: \citeauthor{mo2026lowlight}~et al.\ show that low-light degradation induces hallucinated semantics and unstable scores~\cite{mo2026lowlight}. This observation speaks directly to the problem we target, yet a systematic answer to ``when can it not be used, and what to do then'' is still missing.

\subsection{Explanation Generation and Trustworthiness Evaluation}

Making the model state its rationale is a direct route to deployment trust. \citeauthor{hashimoto2025caption}~et al.\ automatically generate descriptions for surveillance video with a multimodal large language model~\cite{hashimoto2025caption}; \citeauthor{khedher2025trustworthy}~et al.\ constrain VLM--LLM explanations by rules to enhance trustworthiness~\cite{khedher2025trustworthy}; \citeauthor{alradi2026hierarchical}~et al.\ propose a hierarchical vision-language model with comprehensive language descriptions~\cite{alradi2026hierarchical}. On the evaluation side, \citeauthor{du2024causation}~et al.\ provide a comprehensive benchmark for causation understanding of video anomalies~\cite{du2024causation}, and \citeauthor{zhao2025smarthomebench}~et al.\ build a multimodal large language model benchmark for smart-home anomaly detection~\cite{zhao2025smarthomebench}.

Although the generation side is crowded, evaluation of explanation \emph{quality} is nearly empty: existing benchmarks mainly ask whether an anomaly can be detected or described, not whether each claim in the explanation is supported by video evidence and whether the decision survives removing the key evidence. We formalise this gap as computable metrics and a companion protocol.

\subsection{Efficient Deployment, Edge Computing and Privacy}

Real deployments are further constrained by efficiency and compliance. \citeauthor{phan2025aadcnet}~et al.\ propose a multimodal framework for automatic anomaly detection in real-time surveillance~\cite{phan2025aadcnet}; \citeauthor{yun2025edgegnn}~et al.\ move graph-based anomaly detection onto edge devices with efficient adaptive knowledge graph learning~\cite{yun2025edgegnn}. On privacy, \citeauthor{li2026federated}~et al.\ achieve privacy-preserving VAD via federated learning~\cite{li2026federated}. Multi-camera settings add further requirements: \citeauthor{mishra2026dementia}~et al.\ show that large language models improve scene-invariant detection of behaviour of risk across multiple surveillance camera views in dementia residential care~\cite{mishra2026dementia}.

These works reveal a genuine tension between ``larger models'' and ``lower cost''. Our DAA answers it with low-rank adaptation instead of full fine-tuning, while DAG spends computation on genuinely difficult inputs through clip-level dynamic routing and explicit abstention.
